\item A un estudiante se le suministra un gran número de papel de copia, regla, compás, tijeras y una balanza sensible. Corta varias formas en diferentes tamaños, calcula sus áreas, mide sus masas y traza la gráfica de la figura TP1.1. (a) Considere el cuarto punto experimental de la parte superior. ¿Qué tan lejos está verticalmente de la recta de mejor ajuste? Exprese su respuesta como diferencia en la coordenada del eje vertical. (b) Exprese su respuesta como un porcentaje. (c) Calcule la pendiente de la recta. (d) Indique lo que la gráfica muestra, refiriéndose a la forma de la gráfica y los resultados de los incisos (b) y (c). (e) Describa si este resultado debe ser previsto teóricamente. (f) Describa el significado físico de la pendiente.
